\chapter{Introduction}
\label{ch:intro}

% Allow a little extra line stretch so long technical compounds rebreak
% cleanly instead of spilling into the margin (no template packages added).
\setlength{\emergencystretch}{3em}

\section{Introduction}
\label{sec:intro-context}

Over the last two decades almost every activity of economic or social consequence
has moved onto networked software. Commerce, banking, public administration,
healthcare, industrial control and private communication are now mediated by web
applications, mobile clients and the programming interfaces that join them: a
business that once required a physical counter is today a service reachable from
anywhere. That migration has been accompanied by an accumulation of data without
precedent. Identity documents, payment instruments, clinical histories, location
traces and corporate intellectual property all sit in systems that are, by design,
reachable over a public network. Data of this kind is no longer a by-product of
doing business; it is an asset in its own right, and frequently the most valuable
one an organisation holds.

Anything valuable attracts people who want to take it, and their motives are as
varied as the data itself. Some adversaries pursue direct financial gain, through
fraud, payment-card theft or the resale of credentials on criminal markets. Others
pursue extortion, encrypting or exfiltrating an organisation's own data and charging
for its return. Others act on behalf of states, for espionage, sabotage or
geopolitical advantage, targeting infrastructure and industry rather than
individuals. What they share is the route in: a flaw in the software, or in the way
it was configured, that permits something its designer never intended.

Guarding against them is the business of cybersecurity, and doing so credibly takes
more than good intentions at design time. Security is not a property an organisation
can assert about itself; it is one that must be demonstrated against an adversary
actively looking for the exception. This is the purpose of the security audit and,
in particular, of the \emph{penetration test}: a controlled, authorised exercise in
which a professional attacks a system as a real adversary would, to establish
empirically what an attacker could achieve and how mature the defences actually are.
The practice is established enough to be codified in methodologies and standards
such as the PTES, the OWASP Web Security Testing Guide, the OSSTMM and NIST
SP~800-115 \cite{ptes, owasp_wstg, osstmm, nist80015}, and in regulated sectors it
is mandated outright: PCI~DSS requires not only that penetration testing be
performed, but that the weaknesses it finds be corrected \emph{and the testing
repeated to verify the correction} \cite{pcidss}.

That final clause is where the practice meets a problem, and it is the problem this
dissertation addresses. An audit creates value only if its findings are fixed, and a
fix counts only once somebody confirms it. This confirmation step, the
\emph{revalidation} or retest of reported findings, is routine, repetitive and
poorly automated. A security audit ends with a report listing each discovered
vulnerability and the steps needed to reproduce it; once the development team applies
fixes, someone must verify that those fixes actually work, that each finding is
genuinely closed rather than incidentally hidden. The stakes are not academic. Known,
previously identified vulnerabilities that are never properly closed remain a leading
cause of breaches: in a survey of nearly 3{,}000 practitioners, 60\% of breach victims
traced their compromise to a known vulnerability whose patch was available but had not
been applied \cite{ponemon_vulnerability_2019}. Yet this re-testing is still performed
largely by hand, one finding at a time, on every remediation cycle, and it demands
careful judgement to avoid declaring a still-exploitable weakness ``fixed''. It is
exactly the kind of high-volume, judgement-laden, evidence-driven task at which a
well-supervised AI assistant should be able to help. The manual character of
revalidation is examined in detail in Section~\ref{subsec:revalidation}.

Artificial intelligence, meanwhile, is advancing faster than almost any previous computing
technology, and the capabilities of the underlying models are improving on a
timescale measured in months rather than years. Each new generation of large
language model is not only more fluent but more \emph{autonomous}: able to plan,
use tools, write and debug code, and pursue a goal across many steps with little
human direction. Cybersecurity sits squarely in the path of this progress,
because the same reasoning-and-acting ability that makes a model a capable
programming assistant also makes it a capable attacker. Keeping pace with this
acceleration is no longer optional for the defensive side; it is a precondition
for staying relevant.

The offensive use of this technology is already a documented reality rather than
speculation. Purpose-built ``dark'' language models such as WormGPT and FraudGPT,
sold on underground forums as safety-stripped alternatives to mainstream
assistants, lower the barrier to entry for phishing, malware generation and social
engineering \cite{malla2024}. State-affiliated threat actors from several countries
have been observed using mainstream models for reconnaissance, scripting and
translation in support of their operations \cite{openai_threatactors}. Research has
shown that autonomous agents can exploit real vulnerabilities directly from their
public descriptions, with no prior knowledge of the flaw \cite{fang2024oneday}. In
November 2025 this trajectory reached a milestone: the first publicly reported
cyber-espionage campaign in which an AI agent executed the large majority of the
intrusion (reconnaissance, exploitation, credential harvesting and lateral
movement) with only sporadic human oversight \cite{anthropic_espionage}.

The frontier has moved further still. In April 2026 Anthropic previewed
\emph{Mythos}, a model whose cybersecurity capability, described as an emergent
side-effect of broader gains in reasoning and autonomy, proved strong enough that
the company declined to release it openly \cite{anthropic_mythos}. Operating with
little human steering, Mythos discovered thousands of high- and critical-severity
vulnerabilities across widely used software and produced working exploits in hours
for flaws that had survived decades of human review, including long-dormant bugs in
mature, security-focused systems; an independent assessment by the United Kingdom's
AI Security Institute corroborated the scale of these results
\cite{aisi_mythos}. Nor is such capability confined to frontier laboratories: the
autonomous penetration-testing system XBOW climbed to the top of a major public
bug-bounty leaderboard, out-reporting human researchers over a sustained period
\cite{xbow}. The economics are shifting decisively: work that took expert teams
weeks can increasingly be attempted overnight. Figure~\ref{fig:timeline} places
these milestones on a single timeline; barely three years separate the first
commodity malicious chatbot from autonomous exploitation at the frontier.

\begin{figure}[htbp]
\centering
\begin{tikzpicture}[
  >={Stealth[length=2mm]},
  ev/.style={rounded corners=2pt, align=center, inner sep=3pt, text width=2.7cm, font=\scriptsize},
  off/.style={ev, draw=rojo, thick, fill=rojo!6},
  def/.style={ev, draw=azul, thick, fill=azul!6},
]
\draw[->, thick, gris50] (-0.5,0) -- (14.4,0);
\foreach \x/\yr in {1/2023, 5/2024, 9/2025, 13/2026} {
  \draw[gris50, thick] (\x,0.12) -- (\x,-0.12);
  \node[font=\footnotesize\bfseries, text=gris25] at (\x,-0.42) {\yr};
}
\node[off, anchor=south] (worm)   at (1,0.55)  {WormGPT / FraudGPT malicious LLMs};
\node[off, anchor=south] (nation) at (5,0.55)  {Nation-state LLM misuse};
\node[off, anchor=south] (xbow)   at (9,0.55)  {XBOW tops bug-bounty leaderboard};
\node[off, anchor=south] (mythos) at (13,0.55) {\emph{Mythos}: autonomous discovery + exploitation};
\node[def, anchor=north] (bigsleep) at (5,-0.7)  {Big Sleep finds real-world 0-day};
\node[off, anchor=north] (esp)      at (9,-0.7)  {First AI-orchestrated espionage campaign};
\draw[gris50] (1,0)  -- (worm.south);
\draw[gris50] (5,0)  -- (nation.south);
\draw[gris50] (9,0)  -- (xbow.south);
\draw[gris50] (13,0) -- (mythos.south);
\draw[gris50] (5,0)  -- (bigsleep.north);
\draw[gris50] (9,0)  -- (esp.north);
\end{tikzpicture}
\caption[Timeline of the offensive-AI arms race]{Selected milestones in the use of artificial intelligence for cyber-offence (red) and cyber-defence (blue), 2023--2026.}
\label{fig:timeline}
\end{figure}

Defenders are not without their own AI. The same class of models has been turned to
protective ends: Google's Big Sleep agent autonomously found a previously unknown
vulnerability in production software and later helped pre-empt one that attackers
were preparing to use \cite{google_bigsleep}. The difficulty is one of balance.
Offensive capability scales cheaply and runs unattended, whereas much defensive
work remains gated by manual, human-paced procedures. Where the two sides draw on
the same technology, the advantage accrues to whoever removes the human bottleneck
first, and, on the defensive side, many of the most important bottlenecks have
not yet been addressed.

These two threads meet at the revalidation problem. One side of the audit cycle,
the attack, is being automated at remarkable speed; the other, the confirmation that
what the attack found has genuinely been fixed, is still done by hand. This
dissertation proposes to close part of that distance with a tool, named
\emph{revalid}, that applies an AI agent to the revalidation of reported findings:
it reads a penetration-test report, interprets each finding's reproduction steps,
and re-executes them against an authorised target under human supervision, returning
an evidence-backed verdict on whether the finding is closed. The remainder of this
chapter states why the work was undertaken, what it set out to achieve, and how the
document is arranged.

\section{Motivation}
\label{sec:motivation}

Two motivations, one personal and one technical, brought this project about.

The personal one came first. Artificial intelligence applied to cybersecurity is the
area of the discipline I find most compelling, and the one I intend to keep building a
career in. A final-year project is a rare opportunity to work on a problem at that
intersection with enough depth to actually learn the craft rather than sample it, and
to do so across the whole lifecycle of a system: designing it, building it, evaluating
it honestly and writing it up. Choosing a subject that sits squarely between offensive
security and applied AI was therefore a deliberate investment in a direction I want to
keep growing in, and much of what the work taught is reported in
Section~\ref{sec:lessons}.

The technical motivation is the gap the first one led me to. Revalidation is not
unautomated: pentest-management platforms track it as a workflow, and automated
security-validation products genuinely re-run their checks after a fix. But those two
capabilities sit on opposite sides of a line. The platforms that are driven by the
content of a report hand the actual verification to a human, while the products that
automate verification re-run only the findings \emph{they themselves} discovered, from
their own internal model of the target, and cannot take in a report written by an
external auditor. Nothing surveyed in Section~\ref{sec:sota} occupies the intersection:
re-executing the reproduction steps of a specific, externally reported finding and
returning a per-finding verdict on it. That gap is narrow, and closing it is worth
doing precisely because the case it leaves uncovered, an audit performed by someone
else, is the ordinary one for most organisations.

\section{Goal and objectives}
\label{sec:objectives}

The goal of this project is to advance the defensive side of this arms race by
giving security professionals, the cybersecurity architects, auditors and
developers responsible for closing the loop on an audit, an AI-driven tool that
accelerates the revalidation of vulnerabilities through \emph{assistance} rather
than replacement. Concretely, the general objective is:

\begin{quote}
\emph{To automate the retesting of reported vulnerabilities using artificial
intelligence, in order to reduce response time and increase efficiency in security
audits, while keeping a human in control of every action that touches a target.}
\end{quote}

This general objective is refined into four specific objectives, which structure the
work reported in the remaining chapters:

\begin{enumerate}[labelindent=\parindent, leftmargin=*, label=\mbox{\bft{[O\arabic*]}}]
\item \bft{Analyse and structure penetration-test reports}, identifying for each
  finding the elements relevant to its revalidation: description, impact, attack
  vector, affected endpoints and, above all, the ordered steps needed to reproduce
  it.
\item \bft{Interpret reproduction steps into executable actions}, designing an
  AI-based component able to turn the natural-language description of a finding into
  concrete, non-destructive verification commands.
\item \bft{Develop an automated, human-supervised retesting system} that executes
  those verification actions against authorised targets and produces an
  evidence-backed verdict of \emph{still-open}, \emph{fixed} or
  \emph{inconclusive} for each finding.
\item \bft{Evaluate the accuracy and usefulness of the system}, comparing its
  verdicts against ground truth and against the effort of traditional manual
  revalidation, in terms of reliability and time saved.
\end{enumerate}

The shape of that autonomy deserves stating precisely, since O3 is worded around it. The
retest is run by the agent: it reads the finding, composes each verification command,
interprets what comes back and decides what to try next, across as many steps as the
problem takes. What the operator holds is not the work but the \emph{authority} --- each
command is authorised before it touches the target, and the operator may steer the agent,
run a command themselves or end the session at any point.
Section~\ref{subsec:llm-safety} sets out why a wholly unsupervised language-model agent
is the wrong instrument for a verdict that a compliance process will rely on, since
hallucination and over-confident conclusions are not incidental defects of these systems
but documented properties of them, and Section~\ref{sec:design-retest} shows what the
gate buys. The obvious objection is that oversight must cost speed.
Section~\ref{sec:eval-time} measures it, by running the same findings on the same backend
with the gate and without it: the difference was under three minutes in thirty, because
what a retest costs is the model's per-step latency rather than the human round-trip. A
fully unattended mode exists as well and was measured; supervision is the default, not
the only option.

\section{Proposed solution}
\label{sec:solution}

The system developed to meet these objectives, named \emph{revalid}, is a local,
single-operator web application that ingests a penetration-test report, uses a
language model to extract each finding together with its reproduction steps, and
then runs an interactive, sandboxed agent that re-verifies the finding against an
authorised lab target. Three principles distinguish it from the offensive tools
surveyed above, and from the automation reviewed in
Chapter~\ref{ch:sota}. It is \emph{report-driven}: it re-executes the reproduction
steps of a specific, human-reported finding rather than hunting for new
vulnerabilities. It is \emph{human-gated}: the agent proposes each command and a
human operator approves it before it runs, so autonomy is used to accelerate the
expert, never to remove them. And it is \emph{verdict-oriented}: its output is a
conservative, evidence-linked determination that a compliance decision can rely on,
biased toward \emph{inconclusive} whenever the evidence is ambiguous. The precise
gap this design fills is argued in Section~\ref{sec:gap}. Figure~\ref{fig:concept}
shows the pipeline at a glance.

\begin{figure}[htbp]
\centering
\resizebox{\textwidth}{!}{%
\begin{tikzpicture}[
  >={Stealth[length=2.6mm]},
  st/.style={draw=gris50, thick, rounded corners, align=center, text width=2.5cm, minimum height=1.5cm, fill=gris95, font=\scriptsize, inner sep=3pt},
  hl/.style={st, draw=rojo, very thick, fill=rojo!8},
]
\node[st] (rep) at (0,0)       {\bft{Report}\\PDF / JSON};
\node[st] (ext) at (3.2,0)     {\bft{Extract}\\findings + reproduction steps (LLM)};
\node[st] (goal) at (6.4,0)    {\bft{Goal}\\retest objective};
\node[hl] (retest) at (9.6,0)  {\bft{Gated retest}\\agent proposes, operator approves each command (sandbox)};
\node[st] (verdict) at (12.8,0){\bft{Verdict}\\still-open / fixed / inconclusive, with evidence};
\foreach \a/\b in {rep/ext, ext/goal, goal/retest, retest/verdict} {\draw[->, gris50, thick] (\a) -- (\b);}
\end{tikzpicture}%
}
\caption[What \emph{revalid} does]{The \emph{revalid} pipeline, from report to evidence-backed verdict. The gated retest, where human control is exercised, is highlighted.}
\label{fig:concept}
\end{figure}

\section{Development methodology}
\label{sec:methodology}

The project was developed as an individual software-engineering effort following a
lightweight, Kanban-based iterative process, with every functional requirement
tracked as an issue on a project board and integrated through a
continuous-integration gate before merging. Design decisions of any significance
were recorded as architecture decision records, and the software is covered by an
automated test suite across unit, integration and system levels. In keeping with
the pace of the technology it studies, the project itself was built with the
assistance of AI coding tools under human review; the nature and extent of that
assistance is declared, as the regulations require, in Chapter~\ref{ch:methodology}.

\section{Structure of this document}
\label{sec:structure}

The remainder of this dissertation is organised as follows.
Chapter~\ref{ch:sota} sets out the technical background, surveying penetration testing
and the vulnerability-management lifecycle, the successive waves of automation
applied to security testing, the language-model agents that now drive offensive
security, and the extraction of structured information from reports; it then reviews
the state of the art in the tools that revalidate findings today, and closes by
identifying the research gap that \emph{revalid} fills. Chapter~\ref{ch:methodology}
then sets out the methodology followed to build the system: the design-led,
AI-assisted development model, the software process adopted, the technologies used,
and the declaration of AI-tool use required by the regulation.
Chapter~\ref{ch:design} then presents the design of the system itself: the
requirements it answers to, its container and module structure, the containment
topology that enforces its safety constraints, its class and persistence models, how
it reads a report into structured findings, the gated agentic retest loop at its
centre, and the earlier architecture that was built and then deleted.
Chapter~\ref{ch:evaluation} then reports an evaluation of the system
against a real penetration-test report on a live target: the reliability of its
verdicts, and the effort it saves over manual retesting. Chapter~\ref{ch:conclusions},
finally, draws together what the work achieved and what it taught, and sets out the
directions it opens, chief among them a larger-scale evaluation building on the
one reported here. Finally, Appendix~\ref{ch:annex-install} provides a step-by-step
guide to installing and operating the system.
